\section{讨论与局限}
\label{sec:discussion}

初版频率调谐场表示自底向上的对比度，不一定等于分类任务需要的语义重要性，因此“低到高”只表示显著性顺序。第~\ref{sec:future}节简述了用可训练CNN替换路线场的方案，同时保留硬拓扑的明确梯度边界。

几何预处理有意保持简单：热力图平滑一次、Sobel强度沿圈平滑一次、过滤小圈，不做持久性简化。细小树枝可以存在于完整拓扑中，却没有扫描Token；这能节省计算，也可能遗漏有效细节。补边矩形保证等高线树有效，有孔洞的掩膜区域需要Reeb图，平坦区需要确定性的符号破平局规则。

直接读取RGB取消了CNN特征Stem，也同时取消了CNN天然提供的局部纹理聚合。模型只能看到法线和轮廓实际经过的像素，稀疏路线可能漏掉路线之间的证据，法线长度上限也可能降低覆盖率。因此，小图块采样、浅层特征适配器和额外局部扫描都应作为基线或后续扩展，而不是初版中隐藏的组成部分。

带Tag的 $y$ 是紧凑接口，不保证完整保存长分支信息；分裂复制它也不会产生不同的既往历史。终止圈虽然读取两侧，但碗形或非凸区域的射线仍可能重叠。统一Sobel规则还会让强变化区域同时得到更密采样并靠近闭环终点，因此消融实验需要把这两种作用分开。

最后，不规则短序列即使具有线性复杂度，也可能比规则栅格更难高效运行。任何实证比较都需要报告显著性计算、轮廓构建、碰撞查询、空路线比例、填充开销和真实延迟。本文贡献是完整的路线构建和结构分析，而不是精度、鲁棒性、显存效率或吞吐率方面的实测提升；人类阅读过程只是信息调度的启发，不是生物视觉声明。
